\documentclass[10pt,a4paper]{amsart}
\usepackage[margin=19mm]{geometry}
\usepackage{amsmath,amssymb,mathtools}
\usepackage[scr=rsfs,cal=euler]{mathalfa}
\usepackage{xfrac}
\usepackage[unicode,hidelinks,bookmarks=false]{hyperref}
\newcommand{\ie}{i.e.\ }
\newcommand{\cf}{cf.\ }
\newcommand{\resp}{respectively\ }
\newcommand{\Subsection}{\subsection}
\providecommand{\nobreakdash}{\nobreak\hbox{-}}
\setlength{\emergencystretch}{2em}
\multlinegap0pt
\usepackage{bm}
\usepackage{bbm}
\usepackage{dsfont}
\usepackage{xspace}
\usepackage{cancel}
\usepackage{mathtools}
\usepackage{cleveref}
\usepackage{amsthm}
\makeatletter
\@ifundefined{definition}{\newtheorem{definition}{Definition}}{}
\@ifundefined{lem}{\newtheorem{lem}{Lemma}}{}
\@ifundefined{thm}{\newtheorem{thm}{Theorem}}{}
\makeatother
\DeclareMathOperator{\dg}{dg}
\newcommand{\Z}{\mathbb Z}
\newcommand{\aarea}{\mathbf{area}}
\newcommand{\bbounce}{\mathbf{bounce}}
\title[A decomposition of $m$-Dyck paths]{A decomposition of $m$-Dyck paths}
\author{Yuhan Jiang}
\date{Source: arXiv:2508.12521v4 (CC BY 4.0)}
\begin{document}
\maketitle

\noindent\emph{Source and licence.} A decomposition of $m$-Dyck paths. Version arXiv:2508.12521v4 (CC BY 4.0), distributed under the Creative Commons Attribution 4.0 International licence, as recorded in the arXiv licence metadata for this exact version and shown on the abstract page https://arxiv.org/abs/2508.12521v4. Full rights notice: licenses/arxiv-2508.12521-CC-BY-4.0.txt.\par\smallskip

\noindent\textit{Editorial scope.} Counted unit: the Thm whose source label is \texttt{thm:decomp} in the article (source0.tex lines 342--349, source label \texttt{thm:decomp}), delivered together with the article's own complete original proof of it (source0.tex lines 366--388, one \texttt{proof} environment of 23 lines). No mathematics is written, repaired or re-derived: statement and proof are the article's own text line for line. The proof is detached from the statement in the source: the article states the result at lines 342--349 and prints its proof at lines 366--388, and the proof environment's own header names the statement, so the association is the article's own. Required context retained uncounted: lem:hv (lines 283--289); def:decom (lines 304--308). Every definition, display and label of those lines is retained unchanged. Omitted from this package and not redistributed: 1--282, 290--303, 309--341, 350--365, 389--578. Nothing inside a retained excerpt is omitted or altered: every retained body is delivered exactly as the article's lines are written, so no omission receipt of any kind accompanies this package. Citations: the retained excerpts carry no citation command at all, so no bibliography entry of the article is transcribed and no reference is redistributed. Reference adaptation: none was needed and none was made; every reference inside the delivered unit resolves inside it, so no unresolved reference or citation is produced. Printed slips kept literal: no printed word, symbol, hypothesis or number of the counted statement or of its proof was altered. Layout and class: the article's own class is \texttt{amsart} with packages including geometry, amsmath, amsthm, amssymb, comment, thmtools, mathtools, graphicx, enumitem, xcolor, hyperref, cleveref, float, tikz; the wrapper is the lane's pinned \texttt{amsart} editorial wrapper at 10pt on A4, which restates the theorem-like environments the retained text needs and adds the article's own macro definitions that the retained text needs (\texttt{\textbackslash Z}, \texttt{\textbackslash aarea}, \texttt{\textbackslash bbounce}, \texttt{\textbackslash dg}), with extra wrapper packages (bm, bbm, dsfont, xspace, cancel, mathtools, cleveref). The delivered numbering is the unit's own. Assets: no retained span contains a figure environment, a graphics-inclusion command, a PDF-page inclusion, a table or a TikZ picture; no image file is copied into this package, the package declares no asset dependency and no image file is redistributed with it. Licence: version arXiv:2508.12521v4 of this article is distributed under the Creative Commons Attribution 4.0 International licence, as recorded in the arXiv licence metadata for this exact version and shown on the abstract page https://arxiv.org/abs/2508.12521v4. A full rights notice is delivered at licenses/arxiv-2508.12521-CC-BY-4.0.txt. Characters: every retained excerpt is pure ASCII, so the delivered engine, XeLaTeX with its default Latin Modern text font, reports no missing character.
\begin{lem}\label{lem:hv}
For any rational Dyck path $\pi \in L_{mn,n}^+$ for which the bounce path has vertical steps $v_0,v_1,\dots,v_{s}$ and horizontal steps $h_0,h_1,\dots$ (the tail terms are conventionally set to 0), for any $i \in \{0,1,\dots,m-1\}$, and for any $r \geq 0$, we have that
\begin{equation}
\sum_{k = 0}^{r} h_{km+i} = \sum_{j=0}^{rm+i} v_j.
\end{equation}
In particular, $\sum_{k \geq 0} h_{km+i} = n$.
\end{lem}

\begin{definition}\label{def:decom}
Let $\pi \in L_{mn,n}^+$ be an $m$-Dyck path of which the bounce path has horizontal steps $h_0,h_1,\dots$.
For any $i \in \{0,1,\dots,m-1\}$, we define $\pi^i \in L_{n,n}^+$ to be the concatenation of the parts under $\pi$ in the columns specified by $h_{km+i}$ for $k\geq 0$.
In particular, the bounce composition of $\pi^i$ is given by $(h_i,h_{m+i},\dots)$.
\end{definition}

\begin{thm}\label{thm:decomp}
For any $m$-Dyck path $\pi \in L_{mn,n}^+$ and for any $i \in \{0,1,\dots,m-1\}$, let $\pi^i$ be defined as in \cref{def:decom}.
Then, the bounce (resp. area) sequence of $\pi$ is equal to the entrywise sum of the bounce (resp. area) sequences of $\pi^i$, i.e.,
\begin{align}
\aarea(\pi) &= \sum_{i=0}^{m-1} \aarea(\pi^i) \\
\bbounce(\pi) &= \sum_{i=0}^{m-1} \bbounce(\pi^i).
\end{align}
\end{thm}

\begin{proof}[Proof of \cref{thm:decomp}]
By \cref{def:decom}, we see that for any $k \in [mn]$, there exists $i \in \{0,\dots,m-1\}$ and $j \in [n]$ such that $\dg(\pi)'_k = \dg(\pi^i)'_j$.
Conversely, for any $i \in [0,m-1]$ and any $j \in [n]$, there exists $k \in [mn]$ such that $\dg(\pi^i)'_j = \dg(\pi)'_k$.
In other words, $\dg(\pi)$ is a concatenation of the parts of $\pi^i$s.
Then, $\dg(\pi)$ as the conjugate of $\dg(\pi)'$, satisfies
$\dg(\pi)_\ell = |\{k: \dg(\pi)'_k \geq \ell\}| = \sum |\{k: \dg(\pi^i)'_j \geq \ell\}| = \sum_{i=0}^{m-1} \dg(\pi^i)_\ell$.
As the area sequence is determined by the diagram, we have the first equality.

For the second equality, observe that for a weakly increasing $n$-tuple with $c_r$ many entries equal to $r$ for each $r \geq 0$, the $k$-th entry equals $\#\{r \geq 0 : c_0 + c_1 + \cdots + c_r < k\}$.
By \cref{def:decom}, the bounce composition of $\pi^i$ is $(h_i, h_{m+i}, h_{2m+i}, \dots)$, so $\bbounce(\pi^i)$ has $h_{rm+i}$ many entries equal to $r$. Applying the observation and then \cref{lem:hv},
\[
\bbounce(\pi^i)_k
\;=\; \#\Bigl\{r \geq 0 : \sum_{\ell=0}^{r} h_{\ell m + i} < k\Bigr\}
\;=\; \#\Bigl\{r \geq 0 : \sum_{j=0}^{rm+i} v_j < k\Bigr\}.
\]
Summing over $i \in \{0,\dots,m-1\}$ and using that $(i,r) \mapsto j = rm+i$ is a bijection from $\{0,\dots,m-1\} \times \Z_{\geq 0}$ onto $\Z_{\geq 0}$,
\[
\sum_{i=0}^{m-1} \bbounce(\pi^i)_k
\;=\; \#\Bigl\{j \geq 0 : \sum_{\ell=0}^{j} v_\ell < k\Bigr\}
\;=\; \bbounce(\pi)_k,
\]
where the last equality is the same observation applied to $\bbounce(\pi)$, which has $v_j$ many entries equal to $j$.
\end{proof}

\end{document}
